% ===== A1. El paquete 2027 mueve la base en el margen (E6, actualizada) =====
\begin{frame}{Apéndice · El paquete 2027 mueve la base en el margen}
  \vspace{0.2em}
  \begin{columns}[T,onlytextwidth]
    \begin{column}{0.49\textwidth}
      \caja{Amplía la base (en el margen)}{%
        Deducciones de ISR topadas en 96.67\% de los ingresos acumulables, para empresas con ingresos mayores a 50 mdp: un impuesto mínimo de alrededor de 1\% de los ingresos.\par\vspace{0.35em}
        El Senado (27 de octubre) excluyó a las armadoras y ensambladoras de vehículos en depósito fiscal y a los expendedores de gasolina y diésel. La Cámara de Diputados ratificó los cambios el 4 de noviembre.\par\vspace{0.35em}
        Pérdidas fiscales: amortización hasta 50\% de la utilidad fiscal del ejercicio.}
    \end{column}
    \begin{column}{0.49\textwidth}
      \caja{Abarata el incumplimiento}{%
        Programa de regularización con condonación de 100\% de multas, recargos y gastos de ejecución.\par\vspace{0.35em}
        Ampliado a grandes contribuyentes: se eliminó el tope de 300 mdp de ingresos.}
    \end{column}
  \end{columns}
  \vspace{0.6em}
  {\small Las exclusiones reducen el rendimiento de la medida. No cambia la restricción de fondo.\par}
  \vspace{-0.3em}
  \fuente{Congreso de la Unión, minuta de reforma a la LISR 2027; LIF 2027 (transitorio vigésimo primero), aprobada por el Senado el 27 de octubre de 2026. Pérdidas y regularización: minutas aprobadas por la Cámara de Diputados en octubre de 2026.}
\end{frame}

% ===== A2. Cuatro frentes de reforma (E9) =====
\begin{frame}{Apéndice · Cuatro frentes de reforma}
  \vspace{0.3em}
  \begin{columns}[T,onlytextwidth]
    \begin{column}{0.49\textwidth}
      \caja{Reforma fiscal clásica (federal)}{Bases gravables, gastos fiscales, IVA, ISR.}\par\vspace{0.5em}
      \caja{Haciendas locales}{Ingresos propios de estados y municipios (predial, tenencia moderna) y financiamiento de sus funciones.}
    \end{column}
    \begin{column}{0.49\textwidth}
      \caja{Seguridad social}{Pensiones y salud con financiamiento sostenible e inclusión de trabajadores informales.}\par\vspace{0.5em}
      \caja{Empresas del Estado}{Pemex y CFE: reglas claras para los apoyos fiscales y registro transparente de los pasivos contingentes que hoy quedan fuera del presupuesto.}
    \end{column}
  \end{columns}
  \vspace{1.0em}
  {\small Ninguno de los cuatro frentes cierra la cuenta solo.\par}
\end{frame}

% ===== A3. Detalle del techo de endeudamiento =====
\begin{frame}{Apéndice · Detalle del techo de endeudamiento, LIF 2027}
  \vspace{0.4em}
  {\small
  \begin{tabularx}{\textwidth}{@{}>{\raggedright\arraybackslash}p{3.1cm}>{\raggedright\arraybackslash}X>{\raggedright\arraybackslash}X@{}}
    \toprule
    & \textbf{Endeudamiento neto interno} & \textbf{Endeudamiento neto externo}\\
    \midrule
    Gobierno Federal (techo) & Hasta 1.7 billones de pesos & Hasta 13,500 mdd\\\addlinespace
    Pemex & Hasta 160,619.5 mdp & Hasta 5,259.4 mdd\\\addlinespace
    CFE & Hasta 16,742 mdp & Hasta 725 mdd\\
    \bottomrule
  \end{tabularx}}
  \vspace{0.7em}
  \begin{itemize}\setlength{\itemsep}{0.3em}\small
    \item Gobierno Federal: los techos interno y externo son compensables entre sí. Suman alrededor de 1.94 billones de pesos, \dest{4.9\%} del PIB.
    \item El endeudamiento programado es de 1.57 billones de pesos, \dest{4.0\%} del PIB. La meta de RFSP es 3.9\% del PIB.
    \item El techo de Pemex está condicionado a su meta de balance financiero. Los pasivos contingentes no aparecen en el presupuesto ni en los techos de deuda.
  \end{itemize}
  \fuente{LIF 2027, arts. 1o, 2o y 4o; SHCP, CGPE 2027. \% del PIB con el PIB de SHCP.}
\end{frame}
